\documentclass[a4paper,11pt]{article}
\usepackage{lmodern}
\usepackage[T1]{fontenc}
\usepackage[utf8]{inputenc}
\usepackage{amsmath,amssymb,amsthm}
\usepackage{longtable,booktabs,array,calc}
\usepackage{hyperref}
\usepackage[margin=25mm]{geometry}

\newtheorem{theorem}{Theorem}[section]
\newtheorem{proposition}[theorem]{Proposition}
\newtheorem{lemma}[theorem]{Lemma}
\newtheorem{corollary}[theorem]{Corollary}
\newtheorem{definition}[theorem]{Definition}
\newtheorem{remark}[theorem]{Remark}
\newtheorem{observation}[theorem]{Observation}

\providecommand{\tightlist}{\setlength{\itemsep}{0pt}\setlength{\parskip}{0pt}}
\newcounter{none}

\begin{document}

\section{Full Spherical Coverage of Central Projection in Discrete Space
and the Number-Theoretic Necessity of 5-Dimensional Background
Space}\label{full-spherical-coverage-of-central-projection-in-discrete-space-and-the-number-theoretic-necessity-of-5-dimensional-background-space}

\subsection{--- Resolution of the Inverse Mapping Problem via the 5-Axis
10-Face Relay Model
---}\label{resolution-of-the-inverse-mapping-problem-via-the-5-axis-10-face-relay-model}

\textbf{Author:} Noriaki Kihara (WF System Co., Ltd.)

\textbf{Date:} April 2026

\textbf{DOI:}
\href{https://doi.org/10.5281/zenodo.19624957}{10.5281/zenodo.19624957}

\textbf{Prerequisite papers:} -
\href{https://doi.org/10.5281/zenodo.19427780}{Geometric Formulation of
4-Dimensional Space via Central Projection} -
\href{https://doi.org/10.5281/zenodo.19434932}{Geometric Symmetries of
Central Projection: Mathematical Foundations of Multi-Axis Models} -
\href{https://doi.org/10.5281/zenodo.19601592}{Formulation of Structural
Requirements for a Theory of Everything}

\begin{center}\rule{0.5\linewidth}{0.5pt}\end{center}

\subsection{0. Purpose of This Paper}\label{purpose-of-this-paper}

When inversely mapping the entire surface of \(S^n(R)\) (the
\(n\)-dimensional hypersphere of curvature radius \(R\)) onto background
coordinates (tangent hyperplanes) using central projection, two
essential problems arise.

\begin{enumerate}
\def\labelenumi{\arabic{enumi}.}
\tightlist
\item
  \textbf{Equatorial divergence problem}: In regions near the equator as
  seen from the projection center, the metric of the mapping diverges.
\item
  \textbf{South-hemisphere invisibility problem}: A single tangent
  hyperplane can only map one hemisphere of the hypersphere.
\end{enumerate}

Starting from fundamental theorems of number theory, this paper derives
the following:

\begin{itemize}
\tightlist
\item
  The minimum dimension for lattice distances in discrete space to be
  defined without gaps is \textbf{4} (Lagrange's four-square theorem).
\item
  The background space for central projection onto 4-dimensional
  discrete space is \textbf{5-dimensional} (necessity of homogeneous
  coordinates).
\item
  The orthogonal axis structure of 5-dimensional space yields a
  \textbf{5-axis 10-face relay model} that resolves both problems above.
\item
  The 5-dimensional cross-polytope (5-dimensional analogue of the
  regular octahedron) serves as the unit cell for dense packing, and the
  discrete system closes exactly when the curvature radius satisfies
  \(R^2 \in \mathbb{Z}_{>0}\).
\item
  In the subjective space (4-dimensional), this unit cell is mapped as a
  \textbf{4-dimensional hyperrectangular cell}, and positive/negative
  volume (orientation) is naturally defined.
\end{itemize}

All results in this paper are propositions of projective geometry and
discrete geometry, and no physical interpretation is introduced.

\begin{center}\rule{0.5\linewidth}{0.5pt}\end{center}

\subsection{1. Discrete Space and the Sum-of-Squares
Problem}\label{discrete-space-and-the-sum-of-squares-problem}

\subsubsection{1.1 Problem Setting}\label{problem-setting}

In an \(n\)-dimensional discrete space (integer lattice
\(\mathbb{Z}^n\)), the distance \(d\) between lattice points is given by
the generalized Pythagorean theorem:

\[d^2 = a_1^2 + a_2^2 + \cdots + a_n^2, \quad a_i \in \mathbb{Z}\]

This \(d^2\) takes integer values, though \(d\) itself is generally
irrational.

When using discrete space as the foundation for geometry, it is an
essential requirement for isotropic discretization that \textbf{the
square of the distance between any two lattice points can take any
positive integer value}. If some integer \(N\) cannot be expressed as
the sum of \(n\) squares, then no lattice point exists at that distance,
creating \textbf{direction-dependent gaps (anisotropy)} in the discrete
space.

The question is therefore formulated as follows:

\begin{quote}
\textbf{Question}: What is the minimum \(n\) required to express any
positive integer \(N\) as a sum of \(n\) squares of integers:
\(N = a_1^2 + a_2^2 + \cdots + a_n^2\)?
\end{quote}

\subsubsection{\texorpdfstring{1.2 The Case
\(n = 1\)}{1.2 The Case n = 1}}\label{the-case-n-1}

\(N = a_1^2\) holds only when \(N\) itself is a perfect square
(\(1, 4, 9, 16, \ldots\)). The vast majority of positive integers cannot
be expressed as a single square.

\subsubsection{\texorpdfstring{1.3 The Case
\(n = 2\)}{1.3 The Case n = 2}}\label{the-case-n-2}

By Fermat's two-square theorem, the necessary and sufficient condition
for a positive integer \(N\) to be expressible as the sum of two squares
\(N = a_1^2 + a_2^2\) is that every prime factor of \(N\) of the form
\(4k + 3\) appears to an even power.

For example, \(N = 3\) is a prime of the form \(4 \times 0 + 3\), and no
integer solution exists for \(3 = a_1^2 + a_2^2\). Similarly,
\(N = 7, 15, 21, \ldots\) cannot be expressed as the sum of two squares.

Therefore, in a 2-dimensional lattice, no lattice points exist in
directions where the square of the distance takes these values,
resulting in severe anisotropy in the discretization.

\subsubsection{\texorpdfstring{1.4 The Case
\(n = 3\)}{1.4 The Case n = 3}}\label{the-case-n-3}

By Legendre's three-square theorem, a positive integer \(N\)
\textbf{cannot} be expressed as the sum of three squares if and only if
\(N\) is of the form

\[N = 4^a(8b + 7), \quad a \geq 0, \quad b \geq 0\]

Specifically, \(N = 7, 15, 23, 28, 31, 39, 47, 55, 60, \ldots\) cannot
be expressed as the sum of three squares. Infinitely many such values
exist.

Therefore, isotropic discretization is impossible in principle even in a
3-dimensional lattice. No lattice point exists at distances with
\(d^2 = 7\) or \(d^2 = 15\) in certain directions.

\subsubsection{\texorpdfstring{1.5 The Case \(n = 4\): Lagrange's
Four-Square
Theorem}{1.5 The Case n = 4: Lagrange's Four-Square Theorem}}\label{the-case-n-4-lagranges-four-square-theorem}

\textbf{Lagrange's four-square theorem} (1770) states the following:

\begin{quote}
\textbf{Theorem 1.1 (Lagrange).} Every positive integer \(N\) can be
expressed as the sum of four squares of integers. That is, there always
exist \(a_1, a_2, a_3, a_4\) such that
\[N = a_1^2 + a_2^2 + a_3^2 + a_4^2, \quad a_i \in \mathbb{Z}_{\geq 0}\]
\end{quote}

Concrete examples:

\begin{itemize}
\tightlist
\item
  \(1 = 1^2 + 0^2 + 0^2 + 0^2\)
\item
  \(2 = 1^2 + 1^2 + 0^2 + 0^2\)
\item
  \(3 = 1^2 + 1^2 + 1^2 + 0^2\)
\item
  \(7 = 2^2 + 1^2 + 1^2 + 1^2\) (impossible with three squares)
\item
  \(15 = 3^2 + 2^2 + 1^2 + 1^2\) (impossible with three squares)
\end{itemize}

\textbf{4 is the necessary and sufficient minimum dimension.} As shown
in §1.4, \(n = 3\) is insufficient, and for \(n \geq 5\), setting
\(a_5 = 0\) reduces to the \(n = 4\) case, so 4 is sufficient.

\subsubsection{1.6 Relation to Waring's
Problem}\label{relation-to-warings-problem}

Lagrange's theorem corresponds to the \(k = 2\) case of Waring's problem
(1770):

\begin{quote}
\textbf{Waring's problem}: Find the minimum number of terms \(g(k)\)
required to express any positive integer as a sum of \(k\)-th powers.
\end{quote}

{\def\LTcaptype{none} % do not increment counter
\begin{longtable}[]{@{}cc@{}}
\toprule\noalign{}
\(k\) (power) & \(g(k)\) (minimum terms required) \\
\midrule\noalign{}
\endhead
\bottomrule\noalign{}
\endlastfoot
2 & \textbf{4} (Lagrange's theorem) \\
3 & 9 \\
4 & 19 \\
\end{longtable}
}

The fact that \(g(2) = 4\) for squares (\(k = 2\)) means that \textbf{4
is the minimum dimension at which a discrete space can constitute an
isotropic lattice without gaps in distance computation}.

\subsubsection{\texorpdfstring{1.7 The Algebraic Privilege of \(n = 4\):
Quaternions}{1.7 The Algebraic Privilege of n = 4: Quaternions}}\label{the-algebraic-privilege-of-n-4-quaternions}

The algebraic basis for Lagrange's theorem holding at \(n = 4\) lies in
the multiplicativity of the quaternion norm.

The norm of a quaternion \(q = a + bi + cj + dk\)
(\(a, b, c, d \in \mathbb{Z}\)) is defined as

\[N(q) = a^2 + b^2 + c^2 + d^2\]

This norm is multiplicative:

\[N(q_1 q_2) = N(q_1) \cdot N(q_2)\]

This property (\textbf{Euler's four-square identity}) ensures that the
product of two four-square sums is again a four-square sum. This is the
core of the proof of Lagrange's theorem and is a property intrinsic to
the algebraic structure of 4 dimensions (the quaternion algebra
\(\mathbb{H}\)).

Crucially, the quaternions form a \textbf{skew field satisfying the
associative law}:

\begin{itemize}
\tightlist
\item
  \textbf{Real numbers} (1-dimensional): commutative field.
\item
  \textbf{Complex numbers} (2-dimensional): commutative field.
\item
  \textbf{Quaternions} (4-dimensional): non-commutative but satisfying
  the associative law; a skew field.
\item
  \textbf{Octonions} (8-dimensional): non-commutative and \textbf{not
  satisfying the associative law}.
\end{itemize}

The 8-dimensional octonions are optimal in packing efficiency (E8
lattice, Viazovska 2016), but the lack of associativity means that
computational consistency in projection operations (which involve
division) is not guaranteed. \textbf{Since central projection is
essentially a division operation (normalization of homogeneous
coordinates), an algebraic system satisfying the associative law is
required}, thereby excluding 8 dimensions from the background space for
central projection.

\subsubsection{1.8 Summary: Why 4
Dimensions}\label{summary-why-4-dimensions}

Synthesizing the above arguments, three independent grounds require the
dimension of discrete space to be \textbf{4} from number-theoretic
considerations:

\begin{enumerate}
\def\labelenumi{\arabic{enumi}.}
\tightlist
\item
  \textbf{Completeness of distances}: 4 is the minimum dimension at
  which any positive integer can be expressed as a sum of squares
  (Lagrange's theorem). For \(n \leq 3\), certain distances cannot be
  represented.
\item
  \textbf{Algebraic closure}: 4 dimensions possess a field structure
  (quaternions \(\mathbb{H}\)) with multiplicative norms, being the
  unique non-commutative dimension with this property.
\item
  \textbf{Projective consistency}: Central projection (a division
  operation) requires the associative law, and the highest-dimensional
  field satisfying this is the quaternions (4 dimensions).
\end{enumerate}

These do not mean ``4 dimensions happen to be convenient,'' but rather
that \textbf{4 dimensions are uniquely determined by fundamental
theorems of number theory and algebra}.

\begin{center}\rule{0.5\linewidth}{0.5pt}\end{center}

\subsection{2. Central Projection and the Dimension of Background
Space}\label{central-projection-and-the-dimension-of-background-space}

\subsubsection{2.1 Definition of Central Projection
(Recap)}\label{definition-of-central-projection-recap}

The operation of mapping a point \(P\) on \(S^n(R)\) (the
\(n\)-dimensional hypersphere of curvature radius \(R\)) to the point
\(P'\) where the line from the center \(O\) of the hypersphere through
\(P\) intersects the \(n\)-dimensional tangent hyperplane \(T\) is
called central projection.

In the preceding paper {[}1{]}, the pullback metric and Einstein tensor
were derived from central projection onto \(S^4(R)\). The projection
target tangent hyperplane (subjective space) is \textbf{4-dimensional},
and the background space \(S^4(R)\) is a hypersphere in
\(\mathbb{R}^5\).

\subsubsection{2.2 Homogeneous Coordinates and Dimensional
Relationship}\label{homogeneous-coordinates-and-dimensional-relationship}

The projective geometry of \(n\)-dimensional space is described using
\((n+1)\)-dimensional homogeneous coordinates. Points of the projective
space \(\mathbb{P}^n\) are defined as equivalence classes of
\(\mathbb{R}^{n+1} \setminus \{0\}\):

\[[x_1 : x_2 : \cdots : x_{n+1}]\]

Central projection corresponds to dividing the other components by one
component (e.g., \(x_{n+1}\)) of these homogeneous coordinates.

Therefore, \textbf{for an \(n\)-dimensional subjective space, the
background space must be \((n+1)\)-dimensional}, as required by the
structure of projective geometry.

Since §1 derived that the dimension of subjective space is 4, the
dimension of background space is \textbf{5}.

\subsubsection{2.3 Sum of Squares in 5-Dimensional
Space}\label{sum-of-squares-in-5-dimensional-space}

The distance between lattice points in the 5-dimensional background
space \(\mathbb{R}^5\) is given by

\[d^2 = a_1^2 + a_2^2 + a_3^2 + a_4^2 + a_5^2, \quad a_i \in \mathbb{Z}\]

Since Lagrange's theorem showed that any positive integer is
representable at \(n = 4\), the \(n = 5\) case subsumes Lagrange's
theorem by setting \(a_5 = 0\). Moreover, the additional degree of
freedom from five variables increases the number of lattice point
combinations realizing the same \(d^2\), improving lattice isotropy
beyond 4 dimensions.

\subsubsection{2.4 Why No Dimension Beyond 5 Is
Needed}\label{why-no-dimension-beyond-5-is-needed}

As stated in §1.7, 8-dimensional octonions do not satisfy the
associative law. Therefore, central projection from the 5-dimensional
hypersphere \(S^4(R) \subset \mathbb{R}^5\) to a 4-dimensional tangent
hyperplane is completely definable on the algebraic structure of
quaternions (4 dimensions), and consistency is guaranteed.

Even if a background space of 6 or more dimensions were introduced, the
projection target would still be 4-dimensional, and the additional
dimensions would be redundant degrees of freedom not contributing to the
projection operation. \textbf{The dimension of background space is
necessary and sufficient at subjective space dimension + 1 = 5}.

\subsubsection{2.5 Summary: Why 5
Dimensions}\label{summary-why-5-dimensions}

\begin{enumerate}
\def\labelenumi{\arabic{enumi}.}
\tightlist
\item
  \textbf{Structure of projective geometry}: Central projection onto
  \(n\) dimensions requires \((n+1)\)-dimensional homogeneous
  coordinates.
\item
  \textbf{Algebraic constraint}: The projection operations in the
  background space are completed within the field structure of
  quaternions (associative law).
\item
  \textbf{Elimination of redundancy}: 6 or more dimensions are redundant
  degrees of freedom that do not contribute to projection, violating the
  minimality of structural requirements.
\end{enumerate}

\begin{center}\rule{0.5\linewidth}{0.5pt}\end{center}

\subsection{3. Inverse Mapping Problems of Central
Projection}\label{inverse-mapping-problems-of-central-projection}

\subsubsection{3.1 Limitations of a Single Tangent
Hyperplane}\label{limitations-of-a-single-tangent-hyperplane}

Consider central projection onto the tangent hyperplane \(T_N\) at the
north pole \(N\) of \(S^4(R)\). A point \(P\) on \(S^4(R)\) is mapped to
the point \(P'\) where the line from the center \(O\) (center of the
hypersphere) through \(P\) intersects \(T_N\).

This mapping has two essential limitations:

\textbf{Limitation 1 (Equatorial divergence)}: As \(P\) approaches the
equator (the set of points at geodesic distance \(\pi R / 2\) from the
north pole), the angle between line \(OP\) and \(T_N\) approaches 0, so
\(P'\) diverges to infinity on \(T_N\). That is, the metric (Jacobian)
of the mapping diverges near the equator.

\[\lim_{\theta \to \pi/2} |P'| = \infty\]

where \(\theta\) is the colatitude (angle from the north pole) of \(P\).

\textbf{Limitation 2 (South-hemisphere invisibility)}: For points \(P\)
in the south hemisphere beyond the equator, the line \(OP\) does not
intersect \(T_N\) (or intersects on the opposite side, making the
mapping non-injective). Therefore, \textbf{a single tangent hyperplane
can only map one hemisphere of the hypersphere}.

\subsubsection{3.2 The Essence of the
Problem}\label{the-essence-of-the-problem}

These limitations are inherent in the geometric structure of central
projection and cannot be resolved by any modification of the coordinate
system on a single tangent hyperplane.

Symmetry II (equivalence of axes) demonstrated in the preceding paper
{[}2{]} states that ``projection axes can be freely chosen in central
projection.'' That is, constructing a tangent hyperplane at any point
other than the north pole yields an equivalent central projection.

By maximally exploiting this property, full spherical
coverage---impossible with a single tangent hyperplane---is realized
through a \textbf{combination of multiple tangent hyperplanes}. This is
the 5-axis 10-face relay model of the next section.

\begin{center}\rule{0.5\linewidth}{0.5pt}\end{center}

\subsection{4. The 5-Axis 10-Face Relay
Model}\label{the-5-axis-10-face-relay-model}

\subsubsection{4.1 Orthogonal Axis Structure of 5-Dimensional
Space}\label{orthogonal-axis-structure-of-5-dimensional-space}

In 5-dimensional space \(\mathbb{R}^5\), there exist 5 orthogonal axes.
Each axis has positive and negative directions, yielding a total of
\textbf{10 directions}
(\(\pm e_1, \pm e_2, \pm e_3, \pm e_4, \pm e_5\)).

On \(S^4(R) \subset \mathbb{R}^5\), each of these 10 directions
corresponds to a point on the hypersphere. These 10 points are called
\textbf{reference points}, and the tangent hyperplanes at each reference
point are denoted \(T_1^+, T_1^-, T_2^+, T_2^-, \ldots, T_5^+, T_5^-\).

\subsubsection{4.2 Symmetry of the 10-Face
Configuration}\label{symmetry-of-the-10-face-configuration}

The 10 reference points have the following properties on \(S^4(R)\):

\begin{enumerate}
\def\labelenumi{\arabic{enumi}.}
\item
  \textbf{90-degree spacing}: The geodesic distance between any two
  reference points is \(\pi R / 2\) (90 degrees). This follows
  immediately from the fact that the orthogonal axes are at right
  angles.
\item
  \textbf{Equivalence}: The symmetry group of \(S^4(R)\) is \(O(5)\),
  and there exists a rotation mapping any reference point to any other.
  Therefore, all 10 tangent hyperplanes are geometrically equivalent.
\item
  \textbf{Antipodality}: \(T_k^+\) and \(T_k^-\) correspond to the
  positive and negative directions of the same axis and are located at
  antipodal points on the hypersphere.
\end{enumerate}

\subsubsection{4.3 Construction of
Coverage}\label{construction-of-coverage}

The central projection from each reference point \(p_k\) can regularly
map the hemisphere centered at \(p_k\) (the region within geodesic
distance \(\pi R / 2\) from \(p_k\)). Denote this hemisphere as \(H_k\).

\textbf{Proposition 4.1.} The 10 hemispheres
\(H_1^+, H_1^-, H_2^+, H_2^-, \ldots, H_5^+, H_5^-\) completely cover
\(S^4(R)\). That is,

\[S^4(R) = \bigcup_{k=1}^{5} \left( H_k^+ \cup H_k^- \right)\]

\textbf{Proof.} Take any point \(P\) on \(S^4(R)\). Let the
5-dimensional coordinates of \(P\) be \((x_1, x_2, x_3, x_4, x_5)\)
(\(\sum x_i^2 = R^2\)). At least one component satisfies
\(|x_k| \geq R/\sqrt{5}\) (otherwise
\(\sum x_i^2 < 5 \cdot R^2/5 = R^2\), a contradiction).

When \(|x_k| \geq R/\sqrt{5}\), the geodesic distance from \(P\) to the
reference point \((\text{sgn}(x_k)) \cdot R \cdot e_k\) is
\(\arccos(|x_k|/R) \leq \arccos(1/\sqrt{5}) < \pi/2\), so
\(P \in H_k^{\text{sgn}(x_k)}\). \(\square\)

\subsubsection{4.4 The Relay Principle}\label{the-relay-principle}

A point \(P\) on \(S^4(R)\) generally belongs to multiple hemispheres
(the coverage has overlaps). When mapping \(P\), one selects the tangent
hyperplane where \(P\) is mapped most regularly (i.e., with minimum
metric distortion).

Specifically, one computes the geodesic distance \(\theta_k\) between
\(P\) and each reference point \(p_k\), and selects the tangent
hyperplane of the reference point for which \(\theta_k\) is smallest.
This is equivalent to ``choosing the axis corresponding to the
coordinate component with the largest absolute value among \(P\)'s
coordinates.''

\subsubsection{4.5 Transformation Between Tangent
Hyperplanes}\label{transformation-between-tangent-hyperplanes}

All tangent hyperplanes \textbf{share the same center \(O\) (origin)}
and \textbf{have the same curvature radius \(R\)}. Only the \textbf{axis
direction} differs. Therefore, coordinate transformations between
tangent hyperplanes are \textbf{orthogonal transformations}
(combinations of rotations and reflections) in 5-dimensional space, and
are extremely simple.

Specifically, the transformation between the \(k\)-th axis and the
\(l\)-th axis is a 90-degree rotation in the \((k, l)\) plane:

\[\begin{pmatrix} x_k' \\ x_l' \end{pmatrix} = \begin{pmatrix} 0 & -1 \\ 1 & 0 \end{pmatrix} \begin{pmatrix} x_k \\ x_l \end{pmatrix}\]

with all other components unchanged.

\subsubsection{4.6 Resolution of the Equatorial Divergence
Problem}\label{resolution-of-the-equatorial-divergence-problem}

A point \(P\) near the equator as seen from \(T_k\)
(\(\theta_k \to \pi/2\)) satisfies \(\theta_l < \pi/2\) for a different
axis \(l\) (\(l \neq k\)), and is regularly mapped from the tangent
hyperplane \(T_l\).

\textbf{Theorem 4.2.} In the 5-axis 10-face relay model, for any point
\(P\) on \(S^4(R)\), there exists at least one reference point within
geodesic distance \(\arccos(1/\sqrt{5}) \approx 63.4°\) (\(< 90°\)) from
\(P\). Therefore, metric distortion is bounded, and equatorial
divergence does not occur.

\textbf{Proof.} This follows immediately from the estimate
\(|x_k| \geq R/\sqrt{5}\) in the proof of Proposition 4.1. \(\square\)

\subsubsection{4.7 Resolution of the South-Hemisphere
Problem}\label{resolution-of-the-south-hemisphere-problem}

Since the tangent hyperplane \(T_k^-\) corresponding to the antipodal
point \(p_k^-\) exists, the region constituting the ``south hemisphere''
as seen from \(T_k^+\) is regularly mapped from \(T_k^-\). Furthermore,
tangent hyperplanes of other axes also contribute to coverage, so the
combination of all 10 faces covers the entire sphere without exception.

\subsubsection{4.8 Summary}\label{summary}

The 5-axis 10-face relay model maximally exploits the symmetry of
central projection (equivalence of axes) to achieve full spherical
coverage with the following properties:

\begin{enumerate}
\def\labelenumi{\arabic{enumi}.}
\tightlist
\item
  \textbf{Completeness}: The entire surface of \(S^4(R)\) is mapped.
\item
  \textbf{Regularity}: At any point, metric distortion is bounded
  (within \(\arccos(1/\sqrt{5}) \approx 63.4°\)).
\item
  \textbf{Simplicity}: Transformations between tangent hyperplanes are
  90-degree orthogonal transformations, sharing the same origin and
  curvature radius.
\item
  \textbf{Symmetry}: All 10 tangent hyperplanes are geometrically
  equivalent.
\end{enumerate}

\begin{center}\rule{0.5\linewidth}{0.5pt}\end{center}

\subsection{5. Dense Packing by the 5-Dimensional
Cross-Polytope}\label{dense-packing-by-the-5-dimensional-cross-polytope}

\subsubsection{5.1 Definition of the 5-Dimensional
Cross-Polytope}\label{definition-of-the-5-dimensional-cross-polytope}

The \(n\)-dimensional cross-polytope (regular hyperoctahedron)
\(\beta_n\) is the regular polytope in \(\mathbb{R}^n\) whose vertices
are the \(\pm 1\) points on each coordinate axis. The number of vertices
is \(2n\).

In the 5-dimensional case, \(\beta_5\) has the following 10 vertices:

\[(\pm 1, 0, 0, 0, 0), \quad (0, \pm 1, 0, 0, 0), \quad (0, 0, \pm 1, 0, 0), \quad (0, 0, 0, \pm 1, 0), \quad (0, 0, 0, 0, \pm 1)\]

\subsubsection{5.2 Correspondence Between the Cross-Polytope and the
10-Face
Relay}\label{correspondence-between-the-cross-polytope-and-the-10-face-relay}

The 10 reference points constructed in §4 are the points on \(S^4(R)\):

\[\pm R \cdot e_1, \quad \pm R \cdot e_2, \quad \pm R \cdot e_3, \quad \pm R \cdot e_4, \quad \pm R \cdot e_5\]

These are none other than the vertices of the 5-dimensional
cross-polytope \(\beta_5\) scaled by \(R\).

That is, \textbf{the configuration of reference points in the 5-axis
10-face relay model coincides exactly with the vertex structure of the
5-dimensional cross-polytope.} This correspondence is not coincidental
but follows necessarily from the orthogonal axis structure of
5-dimensional space.

\subsubsection{5.3 Partition of Space by the
Cross-Polytope}\label{partition-of-space-by-the-cross-polytope}

The line segments from the center (origin) of the 5-dimensional
cross-polytope to each vertex partition \(\mathbb{R}^5\) into
\(2^5 = 32\) sectors (orthants). On \(S^4(R)\), a corresponding
partition into \(32\) spherical simplices (sections of the hypersphere)
is obtained.

Each sector corresponds to one simplex of the cross-polytope, and
adjacent sectors share faces. This partition is \textbf{isotropic}, with
all sectors being congruent.

\subsubsection{5.4 Condition for Dense
Packing}\label{condition-for-dense-packing}

Consider densely packing \(S^4(R)\) with a discrete lattice. Let the
minimal discrete unit (cell) be the 5-dimensional cross-polytope with
edge length (lattice constant) \(\ell\).

\textbf{Proposition 5.1.} The number \(M\) of cells along a great circle
(1-dimensional geodesic) of \(S^4(R)\) is

\[M = \frac{2\pi R}{\ell}\]

When \(M\) is an integer, the geodesic closes exactly at an integer
multiple of discrete steps.

In the case of 4-axis symmetry, all four axes share the same \(M\), so

\[R = \frac{M \ell}{2\pi}\]

Taking the lattice constant \(\ell\) as the unit of length
(\(\ell = 1\)),

\[R = \frac{M}{2\pi}\]

\subsubsection{\texorpdfstring{5.5 The Integer Condition on
\(R^2\)}{5.5 The Integer Condition on R\^{}2}}\label{the-integer-condition-on-r2}

For the discrete lattice to \textbf{close exactly without phase drift}
on \(S^4(R)\), the number of steps along geodesics must be consistent in
all directions.

The distance \(d\) in the 4-dimensional subjective space is given by the
Pythagorean theorem:

\[d^2 = n_1^2 + n_2^2 + n_3^2 + n_4^2, \quad n_i \in \mathbb{Z}\]

Requiring that the step count \(M\) be an integer when traversing a
great circle, and that distances in all directions be expressible as
sums of squares of integers by Lagrange's theorem, yields the following
condition:

\textbf{Theorem 5.2.} A necessary condition for a 4-dimensional discrete
lattice to close isotropically and exactly on \(S^4(R)\) is that \(R^2\)
(the square of the curvature radius in units of the lattice constant) is
a \textbf{positive integer}.

Under this condition, all steps of the discrete lattice are managed as
integer values, and the phase drift (cumulative error) upon completing
one circuit is zero. \textbf{The curvature radius takes not continuous
real values but discrete values (\(R^2 \in \mathbb{Z}_{>0}\)).}

\begin{center}\rule{0.5\linewidth}{0.5pt}\end{center}

\subsection{6. Mapping Shape in the Subjective Space (4
Dimensions)}\label{mapping-shape-in-the-subjective-space-4-dimensions}

\subsubsection{6.1 Mapping from 5-Dimensional Cross-Polytope to
4-Dimensional Hyperrectangular
Cell}\label{mapping-from-5-dimensional-cross-polytope-to-4-dimensional-hyperrectangular-cell}

When a 5-dimensional cross-polytope cell on \(S^4(R)\) is centrally
projected onto a tangent hyperplane (4-dimensional subjective space),
what shape does it take in the projection target?

The projection of a 5-dimensional cross-polytope into 4 dimensions
depends on the projection direction; for projection along an axis
direction of the cross-polytope (choosing one axis as the projection
direction), the shadow is a \textbf{4-dimensional regular octahedron}
(\(\beta_4\), 8 vertices).

On the other hand, in the relay model, the observer in the subjective
space seamlessly traverses multiple tangent hyperplanes. In this case,
the local discrete lattice is perceived as a lattice with 4-axis
symmetry.

\textbf{Proposition 6.1.} To an observer in the 4-dimensional subjective
space, the minimal cell of the discrete lattice is perceived as a
\textbf{4-dimensional hyperrectangular cell}.

This is for the following reasons:

\begin{enumerate}
\def\labelenumi{\arabic{enumi}.}
\tightlist
\item
  Since the four axes are symmetric, lattice spacings in each axis
  direction are equal.
\item
  Central projection maps straight lines to straight lines (preservation
  of geodesics), so the orthogonal structure of the background lattice
  is preserved as orthogonal structure in the subjective space.
\item
  When a gravitational field (curvature) is present, lattice spacings
  vary with location, but locally the hyperrectangular shape is
  maintained.
\end{enumerate}

\subsubsection{6.2 Maintenance of Density}\label{maintenance-of-density}

When 5-dimensional cross-polytope cells cover \(S^4(R)\) without gaps,
their central projection images---4-dimensional hyperrectangular
cells---also cover the subjective space without gaps.

\begin{itemize}
\tightlist
\item
  \textbf{Near the projection center}: The lattice is observed as an
  orderly orthogonal lattice.
\item
  \textbf{Regions near the equator}: The lattice is radially stretched,
  but relay to another tangent hyperplane prevents excessive distortion.
\item
  \textbf{Boundaries between adjacent tangent hyperplanes}: They
  interweave at 45-degree angles, but the observer in the subjective
  space perceives a continuous lattice.
\end{itemize}

\subsubsection{6.3 Lattice Deformation by Gravitational
Fields}\label{lattice-deformation-by-gravitational-fields}

When a gravitational field exists in the subjective space, the lattice
deforms from a perfect hypercube into a hyperrectangular cell (a
rectangular solid with unequal edge lengths in each direction). However,
two conditions are maintained:

\begin{enumerate}
\def\labelenumi{\arabic{enumi}.}
\tightlist
\item
  \textbf{Topological connectivity}: Adjacent cells share faces, and the
  topology of the lattice is invariant.
\item
  \textbf{Hypervolume preservation}: The 5-dimensional hypervolume of
  each cell (computed in background coordinates as the hypervolume of
  the cross-polytope) is invariant. The ``distortion'' of the lattice
  observed in the subjective space is merely a reflection of the change
  in central projection angle.
\end{enumerate}

\begin{center}\rule{0.5\linewidth}{0.5pt}\end{center}

\subsection{7. Volume Orientation and
Sign}\label{volume-orientation-and-sign}

\subsubsection{7.1 Orientation in 5
Dimensions}\label{orientation-in-5-dimensions}

Hypervolume in \(\mathbb{R}^5\) has a defined \textbf{orientation}. When
the order of the vertices of a 5-dimensional cross-polytope is fixed,
the hypervolume takes a positive or negative value depending on whether
the ordering is right-handed or left-handed.

Mathematically, the sign of the determinant of five edge vectors
\(v_1, v_2, v_3, v_4, v_5\)

\[V = \det(v_1, v_2, v_3, v_4, v_5)\]

determines the orientation, and \(|V|\) gives the hypervolume.

\subsubsection{7.2 Orientation Reversal and Sign of
Volume}\label{orientation-reversal-and-sign-of-volume}

When centrally projecting a cross-polytope cell on \(S^4(R)\), the sign
of the Jacobian of the mapping changes depending on the direction of
projection (positive or negative of the projection axis).

\textbf{Proposition 7.1.} The mapping with reversed orientation of the
5-dimensional cross-polytope (the ``inside-out'' mapping) is naturally
defined as a \textbf{hyperrectangular cell with negative hypervolume} in
the subjective space (4 dimensions).

This is because the extra dimension of 5-dimensional space provides the
distinction between ``front and back,'' enabling a geometric sign
assignment to volume that is impossible in 4-dimensional space alone,
realized through projection from 5 dimensions.

\subsubsection{7.3 Coexistence of Positive and Negative
Volumes}\label{coexistence-of-positive-and-negative-volumes}

Cells with different orientations (positive hypervolume) and (negative
hypervolume) can coexist at the same coordinates in the subjective
space. From the 5-dimensional perspective, these are separate cells
belonging to different sectors, but they are mapped to the same
coordinate region in the 4-dimensional projection.

In this case, the algebraic sum of positive and negative hypervolumes
can be zero. That is, \textbf{through the orientation structure of 5
dimensions, the creation and annihilation of volume are geometrically
expressible}.

\begin{center}\rule{0.5\linewidth}{0.5pt}\end{center}

\subsection{8. Equivalence of Curvature}\label{equivalence-of-curvature}

\subsubsection{8.1 5-Dimensional Curvature and 4-Dimensional Closed
Space}\label{dimensional-curvature-and-4-dimensional-closed-space}

\(S^4(R) \subset \mathbb{R}^5\) is a 4-dimensional hypersurface with
curvature radius \(R\) in 5-dimensional space. For a 4-dimensional
observer living on this hypersurface, the ``curvature in the 5th
dimension'' is not directly observable.

However, such an observer can indirectly perceive the existence of
curvature through the following means:

\begin{enumerate}
\def\labelenumi{\arabic{enumi}.}
\tightlist
\item
  \textbf{Sum of angles in geodesic triangles}: In flat 4-dimensional
  space this equals \(\pi\) (180 degrees), but on \(S^4(R)\) it exceeds
  \(\pi\).
\item
  \textbf{Closure of geodesics}: Traveling straight in any direction,
  one returns to the starting point after a finite distance
  (\(2\pi R\)).
\item
  \textbf{Finiteness of space}: The volume is finite (hypersurface area
  of \(S^4(R) = 8\pi^2 R^4/3\)), but no ``boundary'' exists.
\end{enumerate}

\subsubsection{8.2 Subjective Equivalence}\label{subjective-equivalence}

\textbf{Theorem 8.1.} For an observer in the 4-dimensional subjective
space, the following two descriptions are observationally
indistinguishable (equivalent):

\begin{enumerate}
\def\labelenumi{(\alph{enumi})}
\item
  The subjective space is constructed as central projection onto
  \(S^4(R)\) in \(\mathbb{R}^5\), with curvature radius \(R\) in the
  5th-dimensional direction.
\item
  The subjective space is a closed 4-dimensional space with finite
  volume and no boundary, whose curvature radius computed from the
  deviation of angle sums and closure of geodesics is \(R\).
\end{enumerate}

\textbf{Sketch of proof.} Central projection maps geodesics to straight
lines (preceding paper {[}1{]}), and the sum of angles is determined by
the curvature of \(S^4(R)\). The quantities measurable by the
subjective-space observer (angle sums, geodesic periods, volume) are all
functions of \(R\) alone and contain no additional information about the
structure in the 5th dimension. \(\square\)

This equivalence is the fundamental justification for using central
projection as a geometric framework for gravitational fields. The
5-dimensional structure of the background space is completely encoded as
curvature in the subjective space, and the observer perceives the
background geometry solely through curvature.

\begin{center}\rule{0.5\linewidth}{0.5pt}\end{center}

\subsection{9. Conclusions}\label{conclusions}

\subsubsection{9.1 Results of This Paper}\label{results-of-this-paper}

Starting from fundamental theorems of number theory, this paper has
derived the following:

\begin{enumerate}
\def\labelenumi{\arabic{enumi}.}
\item
  \textbf{Necessity of 4 dimensions}: By Lagrange's four-square theorem,
  the minimum dimension for lattice distances in discrete space to be
  defined isotropically without gaps is 4. This corresponds to the
  algebraic structure of quaternions (multiplicativity of norms,
  associative law).
\item
  \textbf{Necessity of 5 dimensions}: By the structure of projective
  geometry, the background space for central projection onto 4
  dimensions is 5-dimensional. Since octonions (8 dimensions) do not
  satisfy the associative law, they cannot serve as the background space
  for central projection (a division operation).
\item
  \textbf{5-axis 10-face relay model}: The 10 tangent hyperplanes
  corresponding to the positive and negative directions of the 5
  orthogonal axes in 5-dimensional space (= the 10 vertices of the
  5-dimensional cross-polytope) constitute a full spherical coverage of
  \(S^4(R)\). This simultaneously resolves the equatorial divergence
  problem and the south-hemisphere invisibility problem of central
  projection.
\item
  \textbf{Dense packing and the integer condition on \(R^2\)}: A
  discrete lattice with the 5-dimensional cross-polytope as its unit
  cell constitutes an exact closed dense packing on \(S^4(R)\) when the
  curvature radius satisfies \(R^2 \in \mathbb{Z}_{>0}\).
\item
  \textbf{Mapping to 4-dimensional hyperrectangular cells}: The
  5-dimensional cross-polytope cells are mapped as 4-dimensional
  hyperrectangular cells in the subjective space (4 dimensions). The
  orientation structure of 5 dimensions naturally defines positive and
  negative volumes.
\item
  \textbf{Equivalence of curvature}: Curvature in the 5th-dimensional
  direction and the intrinsic curvature of a closed 4-dimensional space
  are indistinguishable to an observer in the subjective space.
\end{enumerate}

\subsubsection{9.2 Scope of This Paper}\label{scope-of-this-paper}

All results in this paper are stated as propositions of projective
geometry, discrete geometry, and number theory, without invoking
physical interpretations (gravitational fields, quantum mechanics,
etc.). The physical meaning of discrete lattices and the correspondence
between hypervolumes of cross-polytopes and physical quantities (energy,
etc.) are left as future work in separate papers.

\begin{center}\rule{0.5\linewidth}{0.5pt}\end{center}

\subsection{References}\label{references}

{[}1{]} N. Kihara, ``Geometric Formulation of 4-Dimensional Space via
Central Projection,'' 2026. DOI:
\href{https://doi.org/10.5281/zenodo.19427780}{10.5281/zenodo.19427780}

{[}2{]} N. Kihara, ``Geometric Symmetries of Central Projection:
Mathematical Foundations of Multi-Axis Models,'' 2026. DOI:
\href{https://doi.org/10.5281/zenodo.19434932}{10.5281/zenodo.19434932}

{[}3{]} N. Kihara, ``Formulation of Structural Requirements for a Theory
of Everything,'' 2026. DOI:
\href{https://doi.org/10.5281/zenodo.19601592}{10.5281/zenodo.19601592}

{[}4{]} J.-L. Lagrange, ``Démonstration d'un théorème d'arithmétique,''
\emph{Nouveaux Mémoires de l'Académie royale des Sciences et
Belles-Lettres de Berlin}, 1770.

{[}5{]} A. Hurwitz, ``Über die Composition der quadratischen Formen von
beliebig vielen Variablen,'' \emph{Nachrichten von der Gesellschaft der
Wissenschaften zu Göttingen}, 309--316 (1898).

{[}6{]} M. Viazovska, ``The sphere packing problem in dimension 8,''
\emph{Annals of Mathematics} \textbf{185}, 991--1015 (2017).

{[}7{]} J. H. Conway and N. J. A. Sloane, \emph{Sphere Packings,
Lattices and Groups}, 3rd ed., Springer (1999).

\end{document}
